% Single paragraph: the class typesets \abstract inside the title block with
% parskip active, so blank lines become visible paragraph spacing on page 1.
% Opening chain, deliberately in three steps: what RSI asks, why the answer turns
% on the training algorithm, and why no benchmark measures that. Reports means
% rather than win rates -- with 10 tasks a win rate moves too much to carry the
% headline, and the mapped scale is comparable across tasks by construction.
\abstract{%
\textbf{Abstract.}~Recursive self-improvement (RSI) asks whether an AI system can improve the process that produces AI systems, so that the next system inherits the improvement. That process is the training algorithm: a better objective or update rule improves the compute\mbox{-}capability exchange rate for every subsequent run, including the one that produces the next agent. Whether RSI is feasible therefore turns on whether an agent can design training algorithms. No benchmark isolates that ability: existing suites are won by collecting data or by tuning hyperparameters, and none tells a change to how a run is executed apart from a change to how the model learns. We present AI4AI\mbox{-}Bench, 10 frozen research repositories spanning 10 training algorithm families. In each task, an agent has 4 hours on one B300 to rewrite the training algorithm; its code is then rerun from scratch for up to 12 hours and scored by a fixed evaluator hidden from the agent, against the repository's original algorithm under the same procedure. Because the 10 metrics are incommensurable, every task is mapped onto one scale on which $0$ is an uninformative model, $0.1$ is the algorithm the repository ships, and $1.0$ is the task optimum. Across 29 configurations of 6 systems on all 10 tasks the mean score is $0.166$, and the best system reaches $0.250$: even the strongest closes under a fifth of the distance between the algorithm that was already there and the optimum. The submissions show where that distance went: most never change how the model learns at all, and the minority that do average $0.226$ against $0.126$ for the rest. More reasoning effort mostly buys the willingness to go there, taking that minority from $8\%$ of submissions to $64\%$ and the mean score from $0.094$ to $0.196$. We release the task suite, the evaluators and every scored submission, so that the measurement can be repeated as these systems change.
}
